% Chapter 12: Observability, Drift Monitoring & Active Learning

\begin{frame}[plain]
    \vfill
    \centering
    {\footnotesize \textbf{\textcolor{deepblue}{\textsc{TaxonVision-MLOps Architecture}}}} \\[0.35cm]
    {\large \textbf{\textcolor{signalorange}{CHAPTER 12}}} \\[0.25cm]
    {\huge \textbf{Observability \& Active Learning}} \\[0.35cm]
    {\normalsize \textcolor{black!70}{Prometheus Telemetry, Concept Drift Detection \& Continuous Human-in-the-Loop Triage}} \\[0.5cm]
    \textcolor{deepblue}{\rule{0.55\linewidth}{1.2pt}} \\[0.5cm]
    \statuspartial
    \vfill
    \script{In Chapter 12, we present our observability and active learning architecture. We must be completely candid about implementation reality: while our core inference pipeline and basic metric endpoints are operational, full drift monitoring, Evidently AI, Prometheus/Grafana scrape daemons, and Label Studio triage loops are high-priority target roadmap items currently in progress or planned for upcoming sprints.}
    \qa{How does MLOps observability differ from traditional software monitoring?}{An ML API can return HTTP 200 OK while failing silently; we monitor average conformal prediction set sizes and drift distances to detect data shifts in real time.}
\end{frame}

\begin{frame}{Observability}{Production Telemetry: Prometheus Metrics \& Grafana (Roadmap)}
    \begin{block}{Domain-Specific Telemetry Beyond HTTP 200}
        \begin{itemize}\small
            \item \textbf{System Metrics:} p95/p99 inference latency ($<50$ ms budget), GPU VRAM utilization \textcolor{signalorange}{\textbf{[Target Metric]}}.
            \item \textbf{Epistemic Metrics:} Average conformal set size ($|C(X)|$), energy-based OOD anomaly rate \textcolor{signalorange}{\textbf{[Target Metric]}}.
            \item \textbf{Operational Metrics:} Triage referral rate to human taxonomists (Label Studio queue depth) \textcolor{deepblue}{\textbf{[Roadmap]}}.
            \item \textbf{Instrumentation Status:} FastAPI \texttt{/metrics} endpoint and basic counters implemented in \texttt{telemetry.py} \textcolor{plantgreen}{\textbf{[Implemented]}}; external Prometheus scraper and Grafana dashboards \textcolor{deepblue}{\textbf{[Target Roadmap]}}.
        \end{itemize}
    \end{block}

    \vspace{0.15cm}
    \statustarget

    \script{\footnotesize Observability in MLOps differs fundamentally from DevOps: an API can return HTTP 200 OK while predicting nonsense. We track conformal set sizes, referral rates, and OOD energy. Transparently, while FastAPI exposes \texttt{/metrics} and core counters in \texttt{telemetry.py}, dedicated Prometheus scrapers and Grafana dashboards are target Phase 3 roadmap infrastructure.}
    \qa{Are Prometheus and Grafana actively scraping the service?}{\footnotesize Not yet. Metric counters exist in \texttt{telemetry.py}, but Prometheus daemons and Grafana dashboards remain on our roadmap.}
    \qa{Why pull-based Prometheus scraping over push?}{\footnotesize Pull scraping prevents telemetry bottlenecks or database timeouts from blocking the inference path.}
\end{frame}

\begin{frame}[fragile]{Telemetry Architecture}{In-Process Monitoring: FastAPI $\rightarrow$ Prometheus $\rightarrow$ Grafana}
    \centering
    \resizebox{0.95\textwidth}{!}{
    \begin{tikzpicture}[
        >=stealth, auto,
        box/.style={rectangle, draw=deepblue, thick, fill=blue!6, text width=2.4cm, align=center, minimum height=1.1cm, rounded corners=3pt, font=\scriptsize},
        db/.style={cylinder, draw=soilbrown, thick, aspect=0.25, fill=orange!10, text width=2.2cm, align=center, minimum height=1.1cm, shape border rotate=90, font=\scriptsize},
        arr/.style={->, thick, >=stealth, draw=black!70}
    ]
        \node[box, draw=plantgreen, fill=green!6] (onnx) {\textbf{1. ONNX Engine}\\\& Conformal Gate\\{\tiny\textcolor{plantgreen}{\textbf{[Active In-Process]}}}};
        \node[box, right=2.4cm of onnx] (fastapi) {\textbf{2. FastAPI Pod}\\\texttt{/metrics} Endpoint\\{\tiny\textcolor{plantgreen}{\textbf{[Active Endpoint]}}}};
        \node[db, right=2.4cm of fastapi] (prom) {\textbf{3. Prometheus}\\\& TSDB Scraper\\{\tiny\textcolor{deepblue}{\textbf{[Target Daemon]}}}};
        \node[box, right=2.4cm of prom] (grafana) {\textbf{4. Grafana}\\Real-Time Dashboards\\{\tiny\textcolor{deepblue}{\textbf{[Target Dashboards]}}}};

        \draw[arr] (onnx) -- node[above, font=\tiny] {In-Process} (fastapi);
        \draw[arr, dashed] (prom) -- node[above, font=\tiny] {Scrape} node[below, font=\fontsize{5pt}{6pt}\selectfont, text=black!60] {15s HTTP Pull} (fastapi);
        \draw[arr, dashed] (prom) -- node[above, font=\tiny] {PromQL} (grafana);
    \end{tikzpicture}
    }

    \vspace{0.3cm}

    \begin{columns}[T]
        \begin{column}{0.31\textwidth}
            \begin{block}{Latency SLA Histograms}
                \tiny
                \textbf{p95 ($<50$\,ms) \& p99 Buckets:}\\
                Tracks pre-processing, ONNX inference execution, and conformal set computation latency without locking.
            \end{block}
        \end{column}
        \begin{column}{0.31\textwidth}
            \begin{block}{Epistemic \& OOD Gauges}
                \tiny
                \textbf{Conformal \& Energy Gauges:}\\
                Real-time prediction set size $|C(x)|$, Helmholtz free energy $E(x)$, and human triage queue referral depth.
            \end{block}
        \end{column}
        \begin{column}{0.31\textwidth}
            \begin{block}{Non-Blocking Pull Scrape}
                \tiny
                \textbf{15s Scrape Architecture:}\\
                FastAPI exposes \texttt{/metrics}; Prometheus pulls asynchronously, ensuring zero inference latency overhead.
            \end{block}
        \end{column}
    \end{columns}

    \script{This diagram illustrates our target telemetry architecture. In the target state, ONNX Runtime passes statistics to FastAPI, Prometheus scrapes them every 15 seconds, and Grafana renders dashboards. Today, the FastAPI /metrics plumbing is in place, while the full Prometheus and Grafana cluster deployment is our next milestone.}
    \qa{Is this telemetry pipeline currently running live in production?}{The FastAPI endpoint and ONNX runtime are active, but the Prometheus collector daemon and Grafana dashboard nodes are target architecture currently being set up.}
    \qa{Why does Prometheus pull rather than the service pushing metrics?}{Pull-based monitoring prevents the application from hanging if the metrics server goes down or becomes slow, ensuring zero impact on the critical inference path.}
\end{frame}

\begin{frame}{Telemetry Architecture}{Target Metric Inventory (Prometheus Roadmap)}
    \begin{block}{Production Metric Scrape Architecture}
        \footnotesize
        Pull-based collection scraping FastAPI \texttt{/metrics} endpoint every 15 seconds.
        \begin{enumerate}
            \item \texttt{taxon\_predictions\_total} \textcolor{plantgreen}{\textbf{[Implemented in \texttt{telemetry.py}]}}
            \item \texttt{taxon\_inference\_latency\_seconds} \textcolor{deepblue}{\textbf{[Target Histogram]}}
            \item \texttt{taxon\_conformal\_set\_size} \textcolor{deepblue}{\textbf{[Target Gauge]}}
            \item \texttt{taxon\_ood\_energy} \textcolor{deepblue}{\textbf{[Target Distribution]}}
            \item \texttt{taxon\_human\_referrals\_active} \textcolor{deepblue}{\textbf{[Roadmap / Target]}}
        \end{enumerate}
    \end{block}

    \script{This slide lists our target metric inventory. Currently in telemetry.py, only the prediction counter is wired up to FastAPI routes. Epistemic metrics like conformal set size and OOD free energy gauges, as well as human referral queue tracking, are defined as constants or design targets awaiting end-to-end integration.}
    \qa{Which metrics are actively collected today?}{Only prediction counts grouped by status in predict.py. Conformal set size, latency histograms, and referral rate metrics are defined as constants but await end-to-end wiring.}
    \qa{Why measure OOD energy instead of standard softmax probabilities?}{Softmax is notoriously overconfident, pushing unknown inputs to high probability classes. Helmholtz free energy provides a more reliable, unnormalized measure of how far an input is from the training manifold.}
\end{frame}

\begin{frame}[fragile]{Drift Monitoring}{Online 1D Marginal Drift: Wasserstein-1 Distance}
    \begin{columns}[T]
        \begin{column}{0.48\textwidth}
            \begin{block}{Wasserstein-1 Distance}
                \vspace{-0.15cm}
                {\footnotesize $$W_1(P, Q) = \int_{-\infty}^{\infty} |F_P(x) - F_Q(x)| \, dx$$}
                \vspace{-0.2cm}
                \begin{itemize}\tiny
                    \item \textbf{Analytical Formulation:} On 1D scalar projections, $W_1(P, Q) = \frac{1}{N}\sum_{i=1}^N |x_{(i)} - y_{(i)}|$ after sorting.
                    \item \textbf{Sorting Complexity:} $\mathcal{O}(N \log N)$ enables real-time stream evaluation over feature projections.
                    \item \textbf{Code Implementation:} Implemented in \texttt{src/taxon\_vision/monitoring/drift.py} for online batch tracking.
                    \item \textbf{Alerting Invariant:} Threshold $\tau_W = 0.15$ triggers triage alerts upon covariate shift from calibration holdout baseline.
                \end{itemize}
            \end{block}
        \end{column}
        \begin{column}{0.48\textwidth}
            \begin{tikzpicture}
            \begin{axis}[
                width=1.02\linewidth,
                height=4.8cm,
                samples=50,
                domain=-3:4,
                ymin=0, ymax=0.48,
                xmin=-3, xmax=4,
                axis lines=left,
                xlabel={\tiny Projected Feature $x$},
                ylabel={\tiny Density $f(x)$},
                xlabel style={font=\tiny, yshift=0.15cm},
                ylabel style={font=\tiny, xshift=-0.15cm},
                tick label style={font=\tiny},
                legend style={font=\fontsize{5pt}{6pt}\selectfont, at={(0.98,0.95)}, anchor=north east},
                legend cell align=left,
                clip=false
            ]
                % Reference P
                \addplot[thick, plantgreen, smooth] {0.39894 * exp(-0.5*x*x)};
                \addlegendentry{Ref Baseline $P$}

                % Shifted Q
                \addplot[thick, signalorange, smooth, dashed] {0.39894 * exp(-0.5*(x - 1.1)*(x - 1.1))};
                \addlegendentry{Shifted Prod $Q$}

                % Wasserstein shift indicator
                \draw[<->, thick, herbired] (axis cs:0, 0.40) -- node[above, font=\tiny\bfseries] {$\Delta \mu = 1.1$} (axis cs:1.1, 0.40);
                \node[anchor=south west, font=\fontsize{5.5pt}{6.5pt}\selectfont, herbired] at (axis cs:0.1, 0.22) {$W_1(P, Q) = 1.10$};
                \node[anchor=north, font=\fontsize{5pt}{6pt}\selectfont, black!70] at (axis cs:0.5, 0.12) {Earth Mover's Distance};
            \end{axis}
            \end{tikzpicture}
        \end{column}
    \end{columns}

    \script{For real-time drift monitoring on streaming inference batches, we employ the 1D Wasserstein-1 Earth Mover's Distance. Because calculating multidimensional drift on every single inference request is computationally prohibitive, we project feature representations onto scalar norms. Sorting these projections allows us to compute Wasserstein distance in $O(N \log N)$ time, flagging distribution shifts in drift.py before accuracy degrades.}
    \qa{Why use Wasserstein distance instead of KL divergence for online drift?}{KL divergence is asymmetric and diverges to infinity when distributions have disjoint support. Wasserstein distance provides a smooth, symmetric, and numerically stable metric reflecting the physical cost of moving mass.}
\end{frame}

\begin{frame}[fragile]{Drift Monitoring}{Offline Multivariate Manifold Drift: RKHS MMD \& Evidently AI}
    \begin{columns}[T]
        \begin{column}{0.48\textwidth}
            \begin{block}{Maximum Mean Discrepancy (RKHS MMD)}
                \vspace{-0.15cm}
                {\scriptsize $$\text{MMD}^2(P, Q) = \mathbb{E}[k(x,x')] - 2\mathbb{E}[k(x,y)] + \mathbb{E}[k(y,y')]$$}
                \vspace{-0.2cm}
                \begin{itemize}\tiny
                    \item \textbf{Kernel Manifold Metric:} Maps embeddings into RKHS $\mathcal{H}$ via Gaussian RBF kernel $k(x,x') = \exp(-\|x-x'\|^2 / 2\sigma^2)$.
                    \item \textbf{Complexity \& Cadence:} $\mathcal{O}(N^2 \cdot D)$ quadratic cost over $D=768/1024$ embeddings; executed as offline scheduled batch job.
                    \item \textbf{Evidently AI Integration:} Target framework generating comprehensive visual drift reports, target drift, and feature correlation degradation.
                    \item \textbf{Retraining Invariant:} Statistically significant MMD divergence ($\alpha < 0.01$) triggers active learning and DVC retrain pipeline.
                \end{itemize}
            \end{block}
        \end{column}
        \begin{column}{0.48\textwidth}
            \centering
            \begin{tikzpicture}[
                scale=0.92,
                refpt/.style={circle, fill=plantgreen, inner sep=1.2pt},
                shiftpt/.style={circle, fill=signalorange, inner sep=1.2pt}
            ]
                % Reference Cluster P
                \draw[thick, plantgreen, fill=plantgreen!10] (-0.9, 0.2) ellipse (1.2cm and 0.8cm);
                \node[font=\tiny\bfseries, plantgreen] at (-0.9, 1.15) {Ref Manifold $P$};
                \node[refpt] at (-1.3, 0.4) {};
                \node[refpt] at (-0.6, 0.0) {};
                \node[refpt] at (-1.1, -0.2) {};
                \node[refpt] at (-0.5, 0.4) {};
                \node[circle, fill=plantgreen!80!black, inner sep=2pt] (muP) at (-0.9, 0.2) {};
                \node[font=\fontsize{5pt}{6pt}\selectfont, below=1pt of muP] {$\mu_P \in \mathcal{H}$};

                % Shifted Cluster Q
                \draw[thick, signalorange, dashed, fill=orange!10] (1.6, -0.3) ellipse (1.2cm and 0.8cm);
                \node[font=\tiny\bfseries, signalorange] at (1.6, 0.65) {Shifted Prod $Q$};
                \node[shiftpt] at (1.2, -0.1) {};
                \node[shiftpt] at (1.8, -0.5) {};
                \node[shiftpt] at (2.0, -0.2) {};
                \node[shiftpt] at (1.4, -0.6) {};
                \node[circle, fill=signalorange!80!black, inner sep=2pt] (muQ) at (1.6, -0.3) {};
                \node[font=\fontsize{5pt}{6pt}\selectfont, below=1pt of muQ] {$\mu_Q \in \mathcal{H}$};

                % MMD Distance Vector
                \draw[->, very thick, herbired] (muP) -- node[above=2pt, font=\fontsize{6pt}{7pt}\selectfont\bfseries] {$\|\mu_P - \mu_Q\|_{\mathcal{H}}$} (muQ);
                \node[anchor=north, font=\fontsize{5.5pt}{6.5pt}\selectfont, herbired] at (0.35, -0.92) {$\text{MMD}(P, Q) = \|\mu_P - \mu_Q\|_{\mathcal{H}}$};

                % Background axes
                \draw[->, thin, black!40] (-2.3, -1.4) -- (3.0, -1.4) node[right, font=\tiny] {$z_1$};
                \draw[->, thin, black!40] (-2.3, -1.4) -- (-2.3, 1.4) node[above, font=\tiny] {$z_2$};
            \end{tikzpicture}
        \end{column}
    \end{columns}

    \script{While 1D Wasserstein evaluates marginal shifts in real time, periodic offline audits require measuring full multivariate manifold drift across high-dimensional feature spaces. Maximum Mean Discrepancy maps embeddings into a Reproducing Kernel Hilbert Space via Gaussian RBF kernels, computing the distance between kernel mean embeddings. In our target architecture, scheduled Evidently AI jobs run MMD to detect subtle covariate and concept shifts, triggering model retraining when divergence is sustained.}
    \qa{Why not compute MMD on every incoming inference request?}{MMD requires computing kernel matrices across sample batches, which has quadratic complexity $O(N^2 \cdot D)$. Running full MMD on every request would degrade serving latency well beyond our fifty-millisecond budget.}
\end{frame}

\begin{frame}{Active Learning}{BADGE Sampling \& Human Review Queue}
    \begin{block}{The Annotation Bottleneck}
        Expert taxonomy time is expensive. Randomly sampling from the unlabelled stream yields redundant, low-diversity annotations.
        \begin{itemize}
            \item \textbf{BADGE Algorithm:} Batch Active learning by Diverse Gradient Embeddings.
            \item \textbf{Dual Criteria:} It computes hypothetical gradient embeddings representing both \textit{epistemic uncertainty} and \textit{feature diversity}.
            \item \textbf{Selection:} Uses $k$-MEANS++ clustering in gradient space to select the optimal batch.
        \end{itemize}
    \end{block}

    \begin{block}{Engineering Rationale}
        By routing only the most mathematically informative samples to Label Studio, we minimise annotation cost while maximising model accuracy improvements per retraining cycle.
    \end{block}

    \vspace{0.2cm}
    \statuspartial

    \script{Annotation by expert taxonomists is expensive and slow. If you send them 10,000 photos of the same common dandelion, you waste their time. BADGE solves this by analyzing the gradients that each unlabelled image would induce in the model. It uses k-MEANS++ clustering in gradient space to pick a batch of images that are both highly uncertain AND visually diverse. This maximizes the model's accuracy improvement per dollar spent on human labeling.}
    \qa{What is the current status of the active learning pipeline?}{The BADGE mathematical sampling algorithm is fully implemented in \texttt{src/taxon\_vision/active\_learning/badge.py}. However, the human review queue integration and automated retraining triggers are in progress.}
\end{frame}

\begin{frame}{Active Learning}{Target Active Learning Loop: Label Studio \& DagsHub (Roadmap)}
    \begin{block}{Architectural \& Operational Target Blueprint}
        \begin{itemize}\small
            \item \textbf{Current Implementation Reality:} Label Studio is \textbf{not yet set up or prepared} in our infrastructure. This slide details the target design for our Phase 3 human-in-the-loop triage engine.
            \item \textbf{Triggering Mechanisms:}
            \begin{itemize}\footnotesize
                \item \textbf{Conformal Ambiguity Filter:} Samples where prediction set size $|\hat{C}(x)| > 1$ or empty sets $|\hat{C}(x)| = 0$.
                \item \textbf{BADGE Diversity Sampling:} Clustering hallucinated gradient embeddings to capture boundary uncertainty \textcolor{plantgreen}{\textbf{[Algorithm Implemented]}}.
            \end{itemize}
            \item \textbf{Annotation Loop (Target):} Automated batch export to Label Studio on DagsHub $\to$ Domain expert review $\to$ DVC dataset commit $\to$ Jenkins retraining webhook \textcolor{deepblue}{\textbf{[Roadmap]}}.
            \item \textbf{Cost Efficiency:} Prevents sending redundant high-confidence common taxa to human reviewers, cutting labeling expenses by $> 70\%$.
        \end{itemize}
    \end{block}

    \vspace{0.15cm}
    \statustarget

    \script{\footnotesize On this slide, we illustrate our target active learning loop. To be completely clear with the team: Label Studio is not yet set up, configured, or integrated into our deployment. The sampling math in BADGE exists in badge.py, but the web UI, annotation ingestion pipelines, and automated DVC retraining webhooks are target roadmap items that our team will build next.}
    \qa{Is Label Studio currently deployed and operational?}{\footnotesize No. There is nothing prepared for Label Studio yet in our infrastructure; it is purely a target architectural blueprint for Phase 3.}
    \qa{Why show this workflow if Label Studio is not yet set up?}{\footnotesize To define the exact data contract and trigger mechanisms so that the engineering team can build the Label Studio connector against a unified specification.}
\end{frame}

\begin{frame}{Active Learning Flow}{Target Closed-Loop Architecture: Stream $\rightarrow$ BADGE $\rightarrow$ Label Studio}
    \begin{center}
    \resizebox{!}{0.48\textheight}{
    \begin{tikzpicture}[
        node distance=0.8cm and 1.2cm,
        stepbox/.style={rectangle, draw=deepblue, thick, fill=blue!5, text width=2.5cm, align=center, minimum height=1.0cm, rounded corners},
        conformalbox/.style={rectangle, draw=deepblue, thick, fill=blue!5, text width=2.8cm, align=center, minimum height=1.0cm, rounded corners},
        arrow/.style={->, thick, >=stealth},
        labelbox/.style={rectangle, draw=signalorange, thick, fill=orange!10, text width=2.5cm, align=center, minimum height=1.0cm, rounded corners}
    ]
        % Nodes
        \node[stepbox] (stream) {\textbf{1. Live Stream} \\ Unlabelled Data \\ {\footnotesize\textcolor{plantgreen}{\textbf{[Active Data]}}}};
        \node[conformalbox, right=of stream] (conformal) {\textbf{2. Conformal}\\\textbf{Filter} \\ Identify ambiguous \\ {\footnotesize\textcolor{plantgreen}{\textbf{[Active Filter]}}}};
        \node[stepbox, right=of conformal] (badge) {\textbf{3. BADGE} \\ Diversity sampling \\ {\footnotesize\textcolor{plantgreen}{\textbf{[Implemented]}}}};
        \node[labelbox, below=of badge] (labelstudio) {\textbf{4. Label Studio} \\ Expert Review \\ {\footnotesize\textcolor{deepblue}{\textbf{[Target System]}}}};
        \node[stepbox, left=of labelstudio] (retrain) {\textbf{5. Retraining} \\ DVC Retrain Trigger \\ {\footnotesize\textcolor{deepblue}{\textbf{[Target Webhook]}}}};

        % Arrows
        \draw[arrow] (stream) -- (conformal);
        \draw[arrow] (conformal) -- (badge);
        \draw[arrow, dashed] (badge) -- (labelstudio);
        \draw[arrow, dashed] (labelstudio) -- (retrain);
        \draw[arrow, dashed] (retrain.south) -- ++(0, -0.55cm) -| (stream.south);
    \end{tikzpicture}
    }
    \end{center}


    \script{\footnotesize This flow shows our target architecture for active learning in production. While streaming observations and conformal filtering are active, and BADGE sampling math is implemented in badge.py, automated Label Studio dispatch and closed-loop retraining webhooks remain roadmap goals.}
    \qa{Which parts of this 5-step loop are implemented today?}{\footnotesize Steps 1 and 2 (streaming data and conformal filtering) and step 3 math (BADGE) are ready. Steps 4 and 5 (Label Studio annotation and automated retraining webhook) are roadmap items.}
    \qa{Why sequence Conformal Filtering before BADGE?}{\footnotesize Conformal filtering quickly isolates ambiguous predictions, shrinking the candidate pool so that BADGE gradient clustering runs in seconds.}
\end{frame}
